\section{The full inverse in coefficient norm}\label{inf:sec-inverse}

Fix any centre order $N$. All constants and exponents in this section are fixed
independently of that order. Only the large-$p$ threshold may
depend on it. We assume the actual centre gap
\eqref{inf:centre-gap}, supplied by Theorem~\ref{inf:integer-gap}.
Write
\[
 \mathcal A=\Hu-R^2,\qquad \mathcal D_0=H_0-R^2,
 \qquad K_0=H_0+R^2.
\]
The actual scalar Dirichlet graph operator is
$J_D=B_{\Phi_N,D}\KD$ and includes the harmonic source layer and
the seed.

\subsection{A detuning-scaled window}

Let $C_{\mathrm{rel}}\ge1$ be a constant in
\eqref{inf:relative-two-norms}. Fix
$E_{\mathrm{inv}}\ge64C_{\mathrm{rel}}$, and define
\begin{equation}\label{inf:inverse-window}
 \Pi=\one_{[-E_{\mathrm{inv}}\tau,E_{\mathrm{inv}}\tau]}
                                    (H_0-R^2),\qquad Q=I-\Pi.
\end{equation}
This is the full ball spectral projection, including every radial
state in the interval. It is distinct from the fixed flow-window
projection. The local count will use
$W_{\mathrm{loc}}=20E_{\mathrm{inv}}$.

\begin{lemma}[Inverse on the full complement]\label{inf:inverse-complement}
For $E_{\mathrm{inv}}\ge64C_{\mathrm{rel}}$ and all $p$ above a
threshold depending on $C_{\mathrm{rel}},E_{\mathrm{inv}},N$,
the operator $\mathcal A_Q=Q\mathcal AQ$ is
invertible on its graph domain, and
\begin{equation}\label{inf:inverse-complement-bound}
 \nrm{K_0\mathcal A_Q^{-1}Q}_{\Mmix,1}
 +\nrm{K_0\mathcal A_Q^{-1}Q}_{\Hc,1}
 \le Cp^2/(E_{\mathrm{inv}}\tau).
\end{equation}
The two inverses coincide on their common sources.
\end{lemma}
\begin{proof}
The projections in \eqref{inf:inverse-window} are diagonal in
angular degree and orthogonal in each radial Hilbert space.
Their norms in both matrix algebras and both underlying spaces
are at most one.
Put $D_Q=Q\mathcal D_0Q$. Radial spectral calculus gives
\begin{equation}\label{inf:ball-complement-bound}
 \nrm{K_0D_Q^{-1}Q}\le1+\frac{2R^2}{E_{\mathrm{inv}}\tau}
\end{equation}
in all four norms.
Since $R/p\le33/16$ for $p\ge64$, its product with
\eqref{inf:relative-two-norms} is at most
\[
 \frac{20C_{\mathrm{rel}}}{p^2}
 +\frac{1089C_{\mathrm{rel}}}{128E_{\mathrm{inv}}}.
\]
The chosen window and $p^2\ge80C_{\mathrm{rel}}$ make this at
most $1/4+1089/8192<1/2$. The window is fixed before this
threshold is imposed.
Factor
\[
 \mathcal A_Q=(I+Q\mathcal VK_0^{-1}K_0D_Q^{-1}Q)D_Q.
\]
The norm-convergent Neumann series in each matrix algebra gives
\[
 \nrm{K_0\mathcal A_Q^{-1}Q}
 \le2\left(1+\frac{2R^2}{E_{\mathrm{inv}}\tau}\right).
\]
Increasing the threshold to include $p^2\ge40E_{\mathrm{inv}}$
absorbs the constant term into $Cp^2/(E_{\mathrm{inv}}\tau)$,
proving \eqref{inf:inverse-complement-bound}.
It maps sources into the graph domain because of its leading
$D_Q^{-1}$ factor and the bound after multiplication by $K_0$.
By the graph realizations of Lemma~\ref{inf:mixed-domain}, the
embedding $\Mmix\subset\Hc$, and uniqueness of the Hilbert
inverse, the two constructions agree on common sources.
\end{proof}

\begin{proposition}[The Schur matrix and its actual gap]
\label{inf:inverse-Schur}
Define the exact Schur matrix at spectral parameter zero by
\begin{equation}\label{inf:effective-matrix}
 \Feff=\Pi\mathcal A\Pi-
               \Pi\mathcal AQ\mathcal A_Q^{-1}Q\mathcal A\Pi.
\end{equation}
For large $p$, each angular index occurs at most once in its index
set. In radial unit vectors with the original angular $G_j$,
its Hilbert inner product has diagonal weights $\eta_j^2$.
It is self-adjoint for that inner product, satisfies
\begin{equation}\label{inf:effective-norm}
 \nrm{\Feff}_{\Mmix,1}+\nrm{\Feff}_{\Hc,1}
                                  \le C_{E_{\mathrm{inv}}}\tau,
\end{equation}
and has the Hilbert gap
\begin{equation}\label{inf:effective-gap}
 \nrm{\Feff v}_\Hc\ge\gamma_p\nrm v_\Hc.
\end{equation}
\end{proposition}
\begin{proof}
The containing fixed window is $20E_{\mathrm{inv}}$.
Lemma~\ref{inf:window-isolation} gives one radial state per angular
index. The normalized radial basis therefore identifies the
mixed norm on this finite space with $\ell^1(2^j)$ exactly.
Self-adjointness of \eqref{inf:effective-matrix} follows from
that of $\mathcal A$ and the complementary Hilbert inverse.
The stated diagonal Hilbert weights follow from
\eqref{inf:true-H-norm}.

On $\Pi$, $K_0$ has norm at most $Cp^2$.
Equation~\eqref{inf:relative-two-norms} gives
$\nrm{\Pi\mathcal V\Pi}\le C\tau$ in both weighted norms.
The feedback factors as
\[
 (\Pi\mathcal VK_0^{-1})Q(K_0\mathcal A_Q^{-1}Q)
                     (\mathcal VK_0^{-1})(K_0\Pi).
\]
Its four factors have sizes $C\tau/p^2$,
$Cp^2/(E_{\mathrm{inv}}\tau)$, $C\tau/p^2$, and $Cp^2$,
respectively. Its norm is at most $C\tau/E_{\mathrm{inv}}$.
The unperturbed diagonal has norm at most $E_{\mathrm{inv}}\tau$.
Submultiplicativity from Lemma~\ref{inf:mixed-algebra} proves
\eqref{inf:effective-norm}, including
exponential angular decay of the feedback.

For $v\in\Pi\Hc$ set
$u=v-\mathcal A_Q^{-1}Q\mathcal A\Pi v$.
It belongs to the graph domain and obeys
$Q\mathcal Au=0$, $\Pi\mathcal Au=\Feff v$.
Orthogonality gives $\nrm u_\Hc\ge\nrm v_\Hc$.
The actual centre-domain gap implies
$\nrm{\mathcal Au}_\Hc\ge\gamma_p\nrm u_\Hc$.
These equalities prove \eqref{inf:effective-gap} for the exact
Schur matrix.
\end{proof}

\begin{corollary}[Finite-matrix coefficient inverse]
\label{inf:effective-inverse}
Let $M=M_{20E_{\mathrm{inv}}}$ be the local-count bound from
Theorem~\ref{inf:local-count}. Then
\begin{equation}\label{inf:effective-inverse-bound}
 \nrm{\Feff^{-1}}_{\ell^1(2^j)}
                         \le C\tau^{-1}p^{5M/3}
\end{equation}
for all sufficiently large $p$.
Its constant and exponent are independent of the centre order.
\end{corollary}
\begin{proof}
Take clusters at separation $D_{\mathrm{sep}}\log p$ in the
actual window. Corollary~\ref{inf:clusters}, applied to its fixed
superset of width $20E_{\mathrm{inv}}$, bounds their cardinalities
by $M$, independently of $D_{\mathrm{sep}}$.
Deleting inter-cluster entries in \eqref{inf:effective-norm}
costs at most $C\tau p^{-D_{\mathrm{sep}}}$ in the coefficient
norm by weighted column sums and in the Hilbert norm by its
weighted row and column sums. Apply
Lemma~\ref{inf:cluster-inverse} with $q_{\mathrm{gap}}=5/3$, $\sigma=1$, and
$c=c_{\mathrm{gap}}$.
Choose once and for all
$D_{\mathrm{sep}}>5M/3+5$.
This proves \eqref{inf:effective-inverse-bound} and its stated
independence.
\end{proof}

\subsection{Recovering the full radial graph}

\begin{proposition}[Mixed-space inverse]\label{inf:mixed-inverse}
The full operator $\mathcal A$ and scalar graph $J_D$ have
inverses on the mixed space, with
\begin{equation}\label{inf:mixed-inverse-bound}
 \nrm{K_0\mathcal A^{-1}}_{\Mmix\to\Mmix}
       +\nrm{J_D^{-1}}_{\Mmix\to\Mmix}
 \le C\tau^{-1}p^{2+5M/3}.
\end{equation}
These inverses agree with their Hilbert realizations.
\end{proposition}
\begin{proof}
For a source $f\in\Mmix$, the exact block solution is
\[
 v=\Feff^{-1}(\Pi f-\Pi\mathcal VQ\mathcal A_Q^{-1}Qf),\qquad
 Qu=\mathcal A_Q^{-1}(Qf-Q\mathcal V\Pi v),\qquad u=v+Qu.
\]
The first off-diagonal factor has norm at most $C/E_{\mathrm{inv}}$,
by \eqref{inf:relative-two-norms} and
\eqref{inf:inverse-complement-bound}. The input $Q\mathcal V\Pi$
has norm at most $C\tau$, and $K_0\Pi$ costs $Cp^2$.
Combining these bounds with \eqref{inf:effective-inverse-bound}
gives $\nrm{K_0u}\le C\tau^{-1}p^{2+5M/3}\nrm f$.
By Lemma~\ref{inf:mixed-domain}, the reconstruction lies in the
mixed graph domain and solves the
full equation, whose Hilbert solution is unique by the centre gap.
This proves the first inverse and identifies it with the Hilbert one.

For the scalar equation $B_{\Phi_N}W=f$, the unitary equation is
$\mathcal A(mW)=-mf$. Equation~\eqref{inf:density-graph-bounds}
bounds multiplication by $m$ and $H_0m^{-1}H_0^{-1}$ on $\Mmix$.
Since $H_0K_0^{-1}$ is an angular-diagonal radial contraction,
it also bounds $H_0m^{-1}K_0^{-1}$ by an absolute constant.
Thus
\[
 J_D^{-1}f=\Delta W
 =(H_0m^{-1}K_0^{-1})(K_0\mathcal A^{-1})(mf),
\]
which proves the second bound in \eqref{inf:mixed-inverse-bound}.
Its signs use $\Delta=-H_0$ and $mW=-\mathcal A^{-1}mf$.
Both inverse identities follow from the graph-domain equation and
Hilbert uniqueness. Every radial and angular source, including the
harmonic layer and the seed, is part of this reconstruction.
\end{proof}

\subsection{Finite head, infinite tail, and membership}

\begin{lemma}[Radial tail and finite-head conversion]
\label{inf:head-tail}
For this lemma the scalar hypotheses on the fixed real finite-centre
chart are $R\in[2p-1,2p+4]$ and
\begin{equation}\label{inf:head-tail-hypotheses}
 \nrm{(P_{\Phi_N}\Delta-\Delta)\KD}_{\Cc\to\Cc}
 +\nrm{(P_{\Phi_N}\Delta-\Delta)\KD}_{\Mmix\to\Mmix}
 \le\frac14.
\end{equation}
There is a fixed $K_{\mathrm{head}}$ independent of the centre order
such that, with coefficient projections
\[
 \Pi_T=\one_{\{j+s\le K_{\mathrm{head}}p\}},\qquad Q_T=I-\Pi_T,
\]
the compressed operator $Q_TJ_DQ_T$ has inverse of norm at most
two in both $Q_T\Cc$ and $Q_T\Mmix$ for large $p$.
Moreover
\begin{equation}\label{inf:head-conversion}
 \nrm{\Pi_Tg}_\Cc\le C_{K_{\mathrm{head}}}\sqrt p\,
                                                   \nrm g_{\Mmix}.
\end{equation}
\end{lemma}
\begin{proof}
The Poisson columns in Lemma~\ref{inf:poisson-columns} have radial
bandwidth one and satisfy the column bound
\eqref{inf:poisson-column-bounds}, including the harmonic column.
On the indicated tail this is at most
$C/((1+2K_{\mathrm{head}})p-1)^2$.
The ratio of neighbouring coefficient weights is at most $e^{2/p}$.
Therefore
\[
 \nrm{R^2Q_T\KD Q_T}_{\Cc\to\Cc}
 \le C/(1+2K_{\mathrm{head}})^2.
\]
For the mixed bound, normalize the radial polynomials by their
squared norms $n_{j,s}=p/(p+2j+2s)$.
The neighbouring ratios of $n_{j,s}$ and their inverses are
bounded by two. After this normalization the tridiagonal
Poisson matrix has row and column sums bounded by the same
quantity up to an absolute factor: adjacent denominators differ
by two and are comparable. Removing entries by $Q_T$ does not
increase those absolute sums. The radial Schur inequality in
each angular block gives
$\nrm{R^2Q_T\KD Q_T}_{\Mmix\to\Mmix}
\le C/(1+2K_{\mathrm{head}})^2$ as well.

Choose the fixed cutoff so that both bounds are below $1/4$.
The hypothesis \eqref{inf:head-tail-hypotheses} bounds the norm of
$Q_T(P_{\Phi_N}\Delta-\Delta)\KD Q_T$ by $1/4$ in both spaces.
Since
\[
 Q_TJ_DQ_T=I_{Q_T}+R^2Q_T\KD Q_T
                   +Q_T(P_{\Phi_N}\Delta-\Delta)\KD Q_T,
\]
a Neumann series proves both tail inverses and their norm bound.
The two series agree on common sources.

On the head, $\wt(2j+2s)\le e^{2K_{\mathrm{head}}}$ and
$n_{j,s}\ge(1+2K_{\mathrm{head}})^{-1}$.
Each angular block contains at most $K_{\mathrm{head}}p+1$ radial
coefficients. Cauchy--Schwarz gives
\[
 \sum_{s\le K_{\mathrm{head}}p-j}\wt(2j+2s)|g_{j,s}|
 \le e^{2K_{\mathrm{head}}}
 \sqrt{(K_{\mathrm{head}}p+1)(1+2K_{\mathrm{head}})}
 \left(\sum_s n_{j,s}|g_{j,s}|^2\right)^{1/2}.
\]
Multiply by $2^j$ and sum in $j$ to prove
\eqref{inf:head-conversion}. The common angular coefficient
structure gives the loss $\sqrt p$.
\end{proof}

\begin{theorem}[Two-sided coefficient inverse]\label{inf:coefficient-inverse}
For every fixed centre order $N$ and all sufficiently large $p$,
$J_D$ is a bounded isomorphism of $\Cc$, and
\begin{equation}\label{inf:coefficient-inverse-bound}
 \nrm{J_D^{-1}}_{\Cc\to\Cc}
 \le C\tau^{-1}p^{5/2+5M/3}
 \le C\frac{p^{b_{\mathrm{inv}}}}{\gamma_p},
 \qquad b_{\mathrm{inv}}=2M+3.
\end{equation}
The constants and exponent are independent of $N$; its threshold
may depend on $N$. The estimates are uniform in
$0<\tau\le20$ and all admissible splits with the stated gap.
\end{theorem}
\begin{proof}
Equation~\eqref{inf:relative-scalar} of
Theorem~\ref{inf:relative-graph} bounds the sum in
\eqref{inf:head-tail-hypotheses} by $C\tau p^{-2}\le1/4$
for $0<\tau\le T$ and $p$ above the fixed-order threshold,
after increasing it depending on the fixed $T$.
Together with $R\in[2p-1,2p+4]$, this supplies the hypotheses
of Lemma~\ref{inf:head-tail}; here $T=20$, and arbitrary fixed
$T$ is covered by Lemma~\ref{inf:fixed-detuning-extension}.
For $f\in\Cc$, first solve $J_Dg=f$ in $\Mmix$ by
Proposition~\ref{inf:mixed-inverse}.
The head $\Pi_Tg$ is finite dimensional and belongs to $\Cc$,
with bound \eqref{inf:head-conversion}.
The full $J_D$ is bounded on $\Cc$ by an absolute constant:
use $J_D=I+R^2\KD+(P_{\Phi_N}\Delta-\Delta)\KD$,
\eqref{inf:poisson-operator-bounds}, $R^2=O(p^2)$, and
\eqref{inf:relative-scalar}.
Thus the tail equation with source
$Q_Tf-Q_TJ_D\Pi_Tg$ has a solution $z\in Q_T\Cc$ and
\[
 \nrm z_\Cc\le C(\nrm f_\Cc+\nrm{\Pi_Tg}_\Cc).
\]
By the embedding into $\Mmix$ and uniqueness of its tail equation,
this $z$ is exactly $Q_Tg$.
This proves membership in $\Cc$ before applying a norm estimate
to the full solution. Hence
\[
 \nrm g_\Cc\le C\nrm f_\Cc+
                      C\sqrt p\nrm g_{\Mmix}
 \le C\tau^{-1}p^{5/2+5M/3}\nrm f_\Cc.
\]
The last bound absorbs the first term since $M\ge1$ and
$\tau\le20$.
The constructed right inverse is injective on $\Cc$ by the
Hilbert realization and its gap, so both inverse identities hold.
Finally
$5/2+5M/3-5/3=5M/3+5/6\le2M+3$.
Using \eqref{inf:centre-gap} gives the second inequality in
\eqref{inf:coefficient-inverse-bound}.
Every constant used in the matrix, reconstruction and tail steps
was fixed before the centre order.
\end{proof}
